\begin{abstract}
A worker lying on the floor after a fall or collapse needs help within minutes, and camera-based alarms for this \pd{} event must run on cheap hardware close to the cameras. Small vision--language models (VLMs) can answer the question ``is a person down?'' without task-specific training, but it is unclear how small, how heavily quantized and how low-resolution they can be before they become unreliable. We study this trade-off with five open VLMs from 450\,M to 4\,B parameters, three quantization levels, three input resolutions, one to eight frames per decision and three prompts, and compare them with a pose-estimation baseline. We build a benchmark from four public fall datasets that separates true \pd{} events from look-alike postures such as kneeling, crouching, sitting and lying on furniture. We measure latency, peak memory and power on an NVIDIA Jetson Orin Nano and a CPU-only device and report the Pareto frontiers. \pending{headline result: smallest configuration within X F1 points of the best; false-alarm rate on hard negatives; streams per device at one decision per second}. A pilot study shows that the best prompt differs between models and that a detailed definition makes the smallest model answer ``yes'' to almost every frame. We release labels, code and per-run results.
\end{abstract}

\begin{IEEEkeywords}
Edge computing, fall detection, model quantization, occupational safety, vision--language models, zero-shot learning.
\end{IEEEkeywords}
